\subsubsection{Décodage des fichiers chromosomiques}\label{sec:lecture}
Lorsque deux versions d'un chromosome diffèrent peu, la durée de lecture peut dépasser celle du calcul des fronts. Le banc isole la conversion du FASTA en deux chaînes de bases. Pour \(B\) octets de fichier et un temps de lecture \(T\), le débit vaut \(B/T\). Le processus relit le même fichier après échauffement. Cette mesure comprend le décodage et les allocations, ainsi que les copies depuis le cache de fichiers de Linux.

Le premier lecteur valide et convertit chaque lettre. Le deuxième ajoute d'un bloc les lignes composées uniquement de lettres A, C, G et T majuscules. Les autres lignes passent par une table de conversion qui conserve les symboles ambigus et ignore les espaces. Le troisième lecteur teste seize caractères à la fois avant de copier la ligne. La figure~\ref{fig:fasta} compare des variantes qui rendent les mêmes caractères sur les six chromosomes.
\chart{cpu/lecture_fasta_chromosomes.png}{Chaque abscisse représente un chromosome. Les compteurs de lecture sur EPYC (a à c) précèdent les rapports de temps appariés et les mesures d'allocation et de mémoire (d à f).}{fig:fasta}{bench/human/chromosome-panel/blocks/index.py}

Pour le chromosome~7, le traitement par ligne ramène les instructions retirées de 25,7 à 5,05 milliards. Les défauts de lecture L1 passent de 16,8 à 16,0 millions. Le lecteur exécute ainsi moins d'instructions, tandis que le nombre absolu de défauts varie peu. Sur le chromosome~1, le traitement explicite par blocs ramène ce compteur de 7,78 à 2,93 milliards sans modifier les 54 demandes d'allocation du processus pendant la lecture.
Le dernier essai réserve la capacité de chaque chaîne dès la lecture de son en-tête. L'index FAI\footnote{Le format \href{https://www.htslib.org/doc/faidx.html}{faidx de HTSlib} donne la longueur, la position et la largeur des lignes de chaque séquence.} fournit cette longueur. Dans neuf blocs distincts, le nombre d'appels à \texttt{new} passe de 53 à 15 pour le chromosome~1. Les octets cumulés demandés passent de 1\,006,64 à 496,79\,Mo. Sur le chromosome~20, le pic RSS moyen passe de 185,89 à 129,00\,Mio. La réservation évite de copier les caractères déjà stockés lorsque la chaîne grandit.

\synthese{Sur le chromosome~1, le lecteur par blocs passe de \(0{,}945\pm0{,}006\) à \(0{,}657\pm0{,}005\)\,s avec l'index. Le gain apparié vaut 1,44 (IC à 95\,\% de 1,43 à 1,45).}
